\section{Conclusion}
\label{sec:conclusion}

This study examines whether structured fine-tuning on Dou Dizhu transfers beyond the source domain, how rule-constrained supervision alters model preferences under explicit prohibitions, and how these behavioral changes are reflected in attention-head-level causal dependence. Across DeepSeek-R1-Distill-Llama-8B, Qwen2.5-7B, and Llama3.1-8B, DD-LoRA exhibits strongly model-dependent transfer on ACL Bench, whereas RC-LoRA achieves higher overall accuracy than DD-LoRA in each case, with category-level effects remaining selective rather than uniform. For Qwen2.5-7B, RC-LoRA produces a substantially larger BCDM than Base, indicating a larger constraint-induced relative preference shift; however, this probability-level change does not translate into improved greedy-decoding compliance, as RC-LoRA selects the prohibited action in all 60 probes. Ablation of individual attention heads further reveals a sparse and heterogeneous redistribution of causal dependence, with a small subset of heads contributing disproportionately to the measured preference shift. Top-$k$ joint ablation produces greater degradation than both global-random and layer-bin-matched random controls, while the non-monotonic effects indicate that head contributions are not simply additive. Overall, these findings support selective cross-task transfer and selective causal reorganization rather than uniform capability gain. They also show that task-transfer performance, constraint-induced preference reallocation, executable rule compliance, and internal causal dependence are distinct outcomes of fine-tuning and should be evaluated separately.
